\documentclass[aspectratio=169,11pt]{beamer}
\usetheme{Madrid}
\usepackage{fontspec}
\setmainfont{LibertinusSerif-Regular.otf}[BoldFont=LibertinusSerif-Bold.otf,ItalicFont=LibertinusSerif-Italic.otf]
\setsansfont{FiraSans-Regular.otf}[BoldFont=FiraSans-Bold.otf,ItalicFont=FiraSans-Italic.otf]
\setmonofont{cmuntt.otf}
\usepackage{polyglossia}
\setdefaultlanguage{russian}
\usepackage{amsmath,booktabs,pgfplots}
\pgfplotsset{compat=1.18}
\definecolor{pine}{HTML}{006C64}
\definecolor{ink}{HTML}{183848}
\definecolor{copper}{HTML}{B75529}
\newcommand{\artm}{\href{https://github.com/bigartm/bigartm}{BigARTM}}
\newcommand{\topicjev}{\href{https://github.com/anatems1/TopicJev/tree/e271339027555160e1a9dc2efa8fc6e71b1a5475}{TopicJev}}
\newcommand{\jev}{\href{https://typesafe.ai/blog/introducing-system-one-models-and-jev}{Jev}}
\newcommand{\laya}{\href{https://huggingface.co/convaiinnovations/laya-typed-decisions}{Laya}}
\newcommand{\luna}{\href{https://developers.openai.com/api/docs/guides/decisions}{OpenAI Decisions}}
\makeatletter\let\rowinput\@@input\makeatother
\newcommand{\basetable}{\begin{tabular}{llrrrr}\toprule Корпус & Метод & PPL $\downarrow$ & NPMI & TD & AMI $\pm s$\\\midrule\rowinput baseline-rows.tex \bottomrule\end{tabular}}
\newcommand{\dectable}{\begin{tabular}{llrrrr}\toprule Корпус & Модель & AMI до & AMI после & Охват & Согласие \\\midrule\rowinput decision-rows.tex \bottomrule\end{tabular}}
\newcommand{\pplplot}{\begin{tikzpicture}\begin{axis}[width=.83\linewidth,height=4.8cm,ybar,bar width=13pt,ymin=0,ymax=110,ytick={0,25,50,75,100},ylabel={\% от PPL униграммы},xtick={1,2,3},xticklabels={PLSA,Сглаживание,Сглаж. + декор.},xmin=.5,xmax=3.5,axis lines*=left,legend columns=2,legend style={at={(.5,1.02)},anchor=south,draw=none},ymajorgrids,grid style={gray!22},tick label style={font=\small}]\addplot[fill=pine,draw=pine] table[x=method,y=ag]{ppl.dat};\addplot[fill=copper,draw=copper] table[x=method,y=ng]{ppl.dat};\legend{AG News,20 Newsgroups}\end{axis}\end{tikzpicture}}

\setbeamertemplate{navigation symbols}{}
\setbeamertemplate{section in toc}[sections numbered]
\AtBeginSection{\title[\insertsectionnumber.\ \insertsectionhead]{Проверка тем документа с помощью моделей принятия решений}}
\setbeamersize{text margin left=8mm,text margin right=8mm}
\newcommand{\point}[1]{\par\medskip{\usebeamercolor[fg]{structure}\bfseries #1}\par\smallskip}
\title[Проверка тематических назначений]{Проверка тем документа с помощью моделей принятия решений}
\subtitle{Jev, TopicJev и эксперимент с BigARTM}
\author{Фанис Хафизов}
\institute[МФТИ]{МФТИ · Вероятностное тематическое моделирование}
\date[26.11.2026]{26 ноября 2026}
\begin{document}
\begin{frame}[plain]\titlepage\note{0–0:30. Представить вопрос доклада: можно ли проверить тему документа по её описанию?}\end{frame}
\begin{frame}{Содержание}
\tableofcontents[hideallsubsections]
\note{0:30–1 мин. Сначала разобрать выбор ответа на примере службы поддержки. Затем определить тематическую смесь и проверить, что даёт повторный выбор темы.}
\end{frame}

\section{Модель Jev}
\begin{frame}{Как Jev обрабатывает обращение в поддержку}
\begin{block}{Текст обращения}
«Я оплатил заказ один раз, но вижу в выписке два списания».
\end{block}
\point{Choice · Куда направить обращение?}
Варианты: платежи, доставка, доступ к аккаунту. \jev{} возвращает выбранный вариант и оценки всех вариантов.
\point{Другие вопросы к тому же тексту}
Бинарный вопрос: «Клиент сообщает о повторном списании?»\\
Оценка по шкале: «Насколько срочно требуется вмешательство?» Для каждого уровня задаётся критерий.
\note{1–3 мин. Попросить слушателей выбрать отдел. По тексту клиент сообщает о двух списаниях, но подтвердить ошибку банка по одной реплике невозможно. Показать, как смысл вопроса определяет смысл ответа. Пример учебный.}
\end{frame}
\begin{frame}{От вопроса к оценкам вариантов}
\point{Вход} Текст обращения, инструкция и описание каждого допустимого ответа.
\point{Вычисление} По описанию TypeSafe, Jev оценивает ответы параллельно. При большом числе вариантов модель сначала оценивает кандидатов по отдельности, затем делает окончательный выбор.
\point{Обучение} TypeSafe называет свой подход RLCD. Цель обучения — согласовать уверенность модели с частотой верных решений.
\note{3–4:30 мин. Одно дело выбрать отдел платежей, другое — оценить надёжность этого выбора. Калибровка проверяется на примерах с известным ответом. Для послойного разбора дальше используем открытую реализацию Laya.}
\end{frame}
\begin{frame}{Открытая архитектура Laya}
\small\texttt{[CLS] вопрос [SEP] [MASK] платежи [MASK] доставка [SEP] текст [SEP]}
\normalsize\point{Контекст} ModernBERT-large кодирует вопрос, варианты и обращение вместе. Каждый маркер получает представление с учётом всего входа.
\point{Голова выбора} К представлениям добавляется вектор типа вопроса. Два слоя Transformer уточняют контекст. Общий скорер выдаёт по одному логиту $z_i$ на маркер.
\point{Нормировка} $p_i=\exp(z_i/T)/\sum_j\exp(z_j/T)$. Оценки зависят от предложенного набора вариантов.
\note{4:30–6 мин. Это открытая архитектура Laya. Возьмите два равных логита: варианты получат по половине. Добавление третьего равного варианта меняет ответ на одну треть.}
\end{frame}
\begin{frame}{Как получить модель выбора из предобученной модели}
\point{Готовый NLI-энкодер} Проверить, следует ли гипотеза «Клиент оспаривает платёж» из текста обращения. Для $K$ вариантов формируем $K$ пар «текст, гипотеза».
\point{Языковая модель} Для однотокенных ответов Yes и No вероятность Yes среди этих двух ответов равна $\sigma(z_{\rm yes}-z_{\rm no})$.
\point{Обучаемая голова выбора} Добавить к энкодеру скорер маркеров. На примерах с правильным индексом $y$ минимизировать $-\log p_y$. Переставлять варианты при обучении и проверять перенос на новые формулировки.
\note{6–8 мин. В первых двух случаях можно начать с готовых весов. В третьем нужны размеченные наборы вариантов. Для вероятности Yes проверить токенизацию в позиции ответа.}
\end{frame}

\section{Постановка задачи тематического моделирования}
\begin{frame}{Тематическое моделирование: дано, найти, критерий}
\point{Дано} Корпус $D$, словарь $W$, частоты слов $n_{dw}$ в документах и число скрытых тем $K$.
\point{Найти} Распределения $\phi_{wt}=p(w\mid t)$ и $\theta_{td}=p(t\mid d)$, для которых
\[
 p(w\mid d)=\sum_{t=1}^{K}\phi_{wt}\theta_{td},\quad
 \phi_{wt},\theta_{td}\ge 0,\quad \sum_w\phi_{wt}=1,\quad \sum_t\theta_{td}=1.
\]
\point{Критерий} В PLSA максимизируем логарифм правдоподобия. В ARTM добавляем регуляризаторы:
\[
 \sum_{d,w}n_{dw}\log p(w\mid d)+\sum_j\tau_j R_j(\Phi,\Theta)
 \longrightarrow \max_{\Phi,\Theta}.
\]
\note{8–10 мин. В PLSA коэффициенты регуляризации равны нулю. Столбец Phi даёт слова темы, столбец Theta даёт смесь тем в документе.}
\end{frame}
\begin{frame}{От тематической смеси к одной рубрике}
\begin{block}{Новость}
«Производитель чипов удвоил прибыль благодаря спросу на ускорители ИИ».
\end{block}
\point{Состав документа} Финансовая и технологическая лексика могут получить вес в одном столбце $\Theta$.
\point{Назначение темы} Исходная модель выбирает $\arg\max_t\theta_{td}$. Можно также описать каждую тему словами из $\Phi$ и попросить модель выбора указать главное содержание новости.
\note{10–11:30 мин. Спросить, какой факт сообщает заголовок. Ответ «финансовые результаты» обоснован ростом прибыли. Теперь обсуждаем, как сравнить такие назначения.}
\end{frame}
\begin{frame}{Три смысла качества}
\point{Предсказание слов · PPL, меньше лучше}
По половине тестового документа оцениваем $\theta_d$ и предсказываем скрытые слова при фиксированной $\Phi$.
\point{Описание темы · NPMI и разнообразие}
NPMI измеряет совместную встречаемость топ-слов. TD@10 — доля различных слов среди $10K$ позиций.
\point{Соответствие рубрикам · AMI}
Сравниваем $\arg\max_t\theta_{td}$ с меткой корпуса. AMI учитывает случайное совпадение и инвариантна к нумерации тем.
\note{11:30–13 мин. Если тема хорошо предсказывает служебные слова, перплексия может улучшиться при слабом описании темы. Корпусные рубрики отражают один принятый способ группировки, документ может затрагивать несколько тем.}
\end{frame}
\section{TopicJev и результаты эксперимента}
\begin{frame}{Что именно проверяет TopicJev}
\point{От модели к описанию} Сначала обучаем \artm{}. Из каждого столбца $\Phi$ берём десять слов и передаём их как описание темы.
\point{От описания к решению} \topicjev{} добавляет Other. В исследованной версии присваиваем тему при $\max_i p_i>2/K$, если победил содержательный вариант.
\point{Модель в эксперименте} Варианты оценивает локальная \laya{}. Для сравнения на тех же документах подготовлен адаптер к \luna{}.
\note{13–15 мин. В демонстрации авторов темы выделял BERTopic. Мы заменили первую ступень на BigARTM. Предобученная модель читает исходный текст и словесные описания. Матрица слов темы после этой стадии не меняется.}
\end{frame}
\begin{frame}{Две коллекции и три конфигурации BigARTM}
\begin{tabular}{lrrr}\toprule Корпус & Обучение & Тест & Темы\\\midrule AG News & 4\,000 & 800 & 4\\20 Newsgroups & 3\,998 & 799 & 20\\\bottomrule\end{tabular}
\point{Одинаковая процедура} Стратифицированные подвыборки исходных train/test. Для 20 NG удалены заголовки, цитаты и подписи. Словарь обучен на train.
\point{18 запусков} PLSA, сглаживание $\Theta$ и $\Phi$, сглаживание с декорреляцией. Для каждой пары «корпус, метод» три инициализации, 25 проходов по коллекции.
\note{15–17 мин. Число тем задано числом рубрик. Метки участвуют в стратификации и в оценивании. Обучение модели и выбор слов для описания проходят без меток документов. После очистки некоторые документы потеряли все или почти все известные слова.}
\end{frame}
\begin{frame}{Сглаживание улучшило предсказание скрытых слов}
\centering\pplplot\par\vspace{1mm}\small 100\% соответствует частотной модели корпуса. Ниже — лучше.\\ Среднее трёх инициализаций.
\note{17–19 мин. Униграммная модель даёт PPL 2172 на AG и 4076 на 20 NG. На двух корпусах сглаживание снижает PPL относительно PLSA. Повторная нормировка к своей униграмме позволяет сравнить относительный эффект при разных словарях.}
\end{frame}
\begin{frame}{Качество тем зависит от выбранной метрики}
\centering\small\basetable\par\medskip\raggedright\normalsize
AMI приведена со стандартным отклонением по трём инициализациям. Выигрыш от добавления декорреляции мал при выбранном весе.
\note{19–21 мин. Рассмотреть 20 NG: сглаживание улучшает PPL и AMI в среднем, но NPMI слегка снижается. При перестановке назначений модуль AMI меньше 0,005, все темы имеют документы. Три запуска показывают чувствительность к старту и не служат доверительным интервалом.}
\end{frame}
\begin{frame}{Какие описания получает верификатор}
\point{20 Newsgroups · тема 3} \textit{space, launch, satellite, nasa, technology, program, lunar, orbit, shuttle, data}
\point{AG News · тема 2} \textit{game, season, iraq, people, team, win, night, killed, said, sunday}
\par\medskip В первом списке угадывается космическая программа. Второй объединяет спортивную лексику и сообщения о войне. Оба списка взяты из измерений со сглаживанием, seed 11.
\note{21–23 мин. Попросить аудиторию сформулировать одно имя для второй темы. Игра, Ирак и погибшие люди могут заставить верификатор выбирать Other. По одним топ-словам это только возможное объяснение, не установленная причина отказов.}
\end{frame}
\begin{frame}{Что изменилось после проверки документов}
\centering\small\dectable\par\medskip\raggedright\normalsize
До и после сравниваются одни и те же 100 документов корпуса. Other входит в расчёт AMI как отдельный кластер.
\point{AG News} Laya назначила содержательную тему 23 документам. На полном наборе AMI снизилась с 0,362 до 0,215.
\note{23–25 мин. В начале выбора 25 новостей из каждой рубрики AG и 5 сообщений из каждой группы 20 NG. В этом режиме Laya ухудшила согласование с внешними рубриками на обеих подвыборках. Утверждение касается этого размера, этих описаний и этого правила Other.}
\end{frame}
\begin{frame}{Почему 77 новостей попали в Other?}
\point{34 из 100} На AG News вариант Other получил наибольшую оценку.
\point{Ещё 43 из 100} Содержательная тема выиграла, но не прошла порог $2/4=0{,}5$. При выборе только по максимуму AMI равна 0,354.
\point{Одна из 43} Для новости о переговорах Fiat и GM максимальная оценка Laya равна 0,241. BigARTM выбрал тему со словами \textit{company, year, million}; после порога ответ — Other.
\note{25–27 мин. Для 20 NG порог равен 0,10, и дополнительного отказа из-за порога не было. Различие между корпусами одновременно связано с текстами, числом тем, их описаниями и калибровкой. Эта диагностика отделяет явный выбор Other от механического действия порога.}
\end{frame}
\begin{frame}{Выводы и проверка с OpenAI Decisions}
\point{BigARTM} Сглаживание снизило PPL на обоих корпусах. Изменения AMI и когерентности оказались меньше и зависели от корпуса и инициализации.
\point{TopicJev + Laya} На парных подвыборках AMI снизилась. Для AG News результат существенно зависел от порога отказа.
\point{OpenAI} Адаптер вызывает \texttt{gpt-6-luna} через \texttt{/v1/decisions}. Зафиксированы те же документы, описания и правило Other. Результаты появляются после платного запуска.
\note{27–30 мин. Отдельно оценивать калибровку на валидационной выборке и качество при одинаковом охвате. Чтобы проверить соответствие именно текстовому описанию, нужна независимая разметка пар документ–тема. Рубрики корпуса отвечают на другой вопрос.}
\end{frame}
\end{document}
